\section{Anexo - Retrospectiva general y cierre documental}

El presente anexo tiene como objetivo realizar una retrospectiva panorámica del Trabajo Profesional, evaluando no solo el alcance funcional del Producto Mínimo Viable (MVP) entregado, sino también el inmenso ecosistema documental, organizativo y de control de calidad construido alrededor de \textbf{SocioUnido}. Esta visión holística permite dimensionar el verdadero impacto y profesionalismo del proyecto.

\subsection{Identidad comercial y validación}

Desde su concepción, SocioUnido fue tratado como un producto real destinado a competir en el mercado de la gestión deportiva. Para respaldar esta visión, se consolidaron los siguientes canales de comunicación e identidad:

\begin{itemize}
    \item \textbf{Canal de YouTube:} Se creó el canal oficial de la plataforma (\url{https://www.youtube.com/@SocioUnido}), el cual funciona como el centro de medios principal. Actualmente aloja los videos de validación del producto con clientes reales, las presentaciones académicas intermedias y la totalidad de los videotutoriales de uso de la plataforma.
    \item \textbf{Página publicitaria (Landing Page):} Pila fundamental para nuestra imagen como marca, se desarrolló e implementó una web comercial oficial con dominio propio, diseñada para la captación y conversión de clientes (\url{\textcolor{blue}{URL A CONFIRMAR}}).
\end{itemize}

El enfoque comercial de SocioUnido no es teórico; fue validado de forma empírica y constante gracias a las iteraciones guiadas por nuestros tutores, y posteriormente refinado mediante reuniones de validación con trabajadores de clubes reales y encuestas de usabilidad a socios y usuarios finales.

\subsection{Ecosistema GitHub: Organización y repositorios}

Absolutamente toda la implementación técnica, lógica y documental del Trabajo Profesional reside de forma centralizada en nuestra organización oficial de GitHub: \url{https://github.com/TrabajoProfesional-AGGZ}. 

La magnitud del proyecto exigió una arquitectura sumamente distribuida, componiéndose de \textbf{27 repositorios y 2 proyectos Kanban}.

\subsubsection*{Gestión ágil y control de incidencias (Proyectos Kanban)}
\begin{itemize}
    \item \textbf{Centro de monitoreo de tareas:} Tablero estructurado bajo la metodología Kanban que rigió la ejecución de los 13 \textit{sprints} de desarrollo. Permitió gestionar el \textit{backlog} de manera dinámica, asignando responsables, pesos (esfuerzo) y prioridades a cada historia de usuario.
    \item \textbf{Centro de monitoreo de bugs:} Espacio homólogo dedicado exclusivamente a la detección, seguimiento y resolución de errores (\textit{bugs}) y deudas técnicas reportadas por los auditores automáticos (SonarCloud/OpenClaw).
\end{itemize}

\subsubsection*{Repositorios privados (Lógica de negocio e infraestructura)}
Con el fin de resguardar la propiedad intelectual, la seguridad criptográfica y los derechos de autor de las implementaciones subyacentes, los repositorios \textit{back-end} se mantienen en carácter privado. Sin embargo, \textbf{ante cualquier requerimiento académico, de evaluación o interés técnico justificado, el equipo está totalmente a disposición para otorgar los accesos correspondientes} (contactando a los correos que figuran en la carátula de este documento).

Cabe destacar que el código fuente de estos repositorios está documentado internamente siguiendo los más altos estándares y buenas prácticas (\textit{self-documenting code}), garantizando que futuros desarrolladores puedan comprender la lógica del sistema de forma clara y completa.

\begin{itemize}
    \item \texttt{microservicio-club}: Gestiona la lógica central (socios, disciplinas, reservas y eventos).
    \item \texttt{microservicio-pagos}: Orquesta las transacciones financieras y pasarelas externas.
    \item \texttt{microservicio-autenticacion}: Administra la identidad, seguridad y roles de los usuarios.
    \item \texttt{microservicio-acceso}: Controla la validación criptográfica (TOTP) y lectura de QRs.
    \item \texttt{microservicio-analiticas}: Calcula y distribuye métricas y proyecciones de morosidad.
    \item \texttt{microservicio-bot-conversacional}: Contiene la lógica asíncrona del bot de Telegram.
    \item \texttt{gateway}: \textit{API Gateway} que unifica y protege el acceso a la red interna.
    \item \texttt{healthchecker}: Monitor de disponibilidad global del ecosistema y de los clubes clientes.
    \item \texttt{openclaw}: Agente de QA especializado para testing intensivo de ciberseguridad.
    \item \texttt{toolkit-financiero}: Herramientas de cálculo, proyección y análisis comercial de la solución.
\end{itemize}

\subsubsection*{Repositorios públicos (Interfaces y portales documentales)}
La totalidad de las interfaces cliente y la documentación asociada a los repositorios privados se mantienen en carácter público para demostrar la transparencia, alcance y arquitectura del proyecto sin comprometer credenciales sensibles.

\textbf{Interfaces gráficas (\textit{Front-end}):} \\
Dentro de estos repositorios reside no solo el código, sino también su propia página documental (JustTheDocs) que expone las vistas y maquetas del sistema.
\begin{itemize}
    \item \texttt{plataforma-web}: Panel de administración (\textit{dashboard}) para los directivos del club.
    \item \texttt{aplicacion}: Aplicación móvil (PWA) interactiva para los socios.
    \item \texttt{control-de-accesos}: Aplicación móvil (PWA) para el escaneo de QRs por parte de los empleados.
\end{itemize}

\textbf{Portales documentales de la arquitectura (\textit{Back-end}):} \\
Para cada microservicio privado existe un repositorio espejo terminado en \texttt{-doc}. Estos alojan portales generados con JustTheDocs que explican exhaustivamente la arquitectura del servicio original: diagramas de contexto, justificación de \textit{stack} tecnológico, métricas automáticas de implementación (actividad de \textit{commits}, lenguajes, contribuciones) y el listado de \textit{endpoints} públicos para integraciones externas.
\begin{itemize}
    \item \texttt{microservicio-club-doc}
    \item \texttt{microservicio-pagos-doc}
    \item \texttt{microservicio-autenticacion-doc}
    \item \texttt{microservicio-acceso-doc}
    \item \texttt{microservicio-analiticas-doc}
    \item \texttt{microservicio-bot-conversacional-doc}
    \item \texttt{gateway-doc}
    \item \texttt{healthchecker-doc}
\end{itemize}

\subsection{El inmenso valor documental del proyecto}

El hito más destacable de SocioUnido, a la par de su impresionante alcance tecnológico y su sólido rendimiento bajo el marco \textit{Scrum} (ver Anexo - Registro de horas PPS), es su aporte documental. Se han construido \textbf{16 páginas plenamente documentales con JustTheDocs} que abarcan el Trabajo Profesional de punta a punta:

\begin{itemize}
    \item \textbf{8 páginas} detallando la arquitectura de los microservicios, \textit{gateway} y monitor de salud.
    \item \textbf{3 páginas} exponiendo las interfaces visuales (Plataforma web y PWAs).
    \item \textbf{5 páginas de gestión global:}
    \begin{itemize}
        \item \texttt{.github}: Centro neurálgico que conecta todos los repositorios, resumiendo el producto y presentando al equipo.
        \item \texttt{scrum}: Documentación profunda de las ceremonias, métricas y evolución de los 13 \textit{sprints} de desarrollo.
        \item \texttt{bitacora}: Reportes detallados del avance a lo largo de las 37 semanas del proyecto.
        \item \texttt{ideas-productos}: El proceso creativo y analítico (5Cs, 4Ps) de las 13 ideas primigenias evaluadas.
        \item \texttt{documentacion}: El repositorio maestro que compila todos los entregables académicos.
    \end{itemize}
\end{itemize}

Particularmente, el repositorio \texttt{documentacion} aloja una vasta cantidad de material complementario generado durante estos meses de trabajo:
\begin{itemize}
    \item El Anteproyecto y las versiones iterativas de la Entrega Intermedia (junto con su video).
    \item Evidencia de la validación del producto (entrevistas a trabajadores y encuestas de usabilidad).
    \item El Informe Final definitivo.
    \item Múltiples redacciones de discursos comerciales (\textit{Pitch / Elevator Pitch}) adaptados a clientes, inversores y académicos.
    \item Desarrollo de la identidad visual: Marca, logotipos, íconos (\textit{marca blanca} y simulando clubes reales).
    \item Diagramas del Modelo C4 (arquitectura de sistemas) y diagramas DER (capa de datos).
    \item Más de \textbf{25 anexos} formales analizando múltiples aristas técnicas, legales y de gestión del sistema.
    \item Decenas de videotutoriales explicativos de la plataforma y las aplicaciones, acompañados por sus respectivos manuales de uso en formato texto.
\end{itemize}

\subsection{Conclusión de la entrega}

El código fuente en su totalidad cuenta con más de \textbf{1.200} \textit{commits} registrados y \textbf{290} revisiones cruzadas (\textit{Pull Requests}) evaluadas. El volumen estructural desarrollado impactó en aproximadamente \textbf{400.000} líneas de código a lo largo de este año de desarrollo continuo.

Por esto, y por el impacto real del producto construido, el equipo no podría estar más orgulloso de los resultados obtenidos. Tenemos la profunda convicción de que el Trabajo Profesional que estamos entregando representa el más alto estándar de excelencia institucional. Haber fortalecido el proyecto en todos los aspectos evaluables (tecnológicos, comerciales, documentales y arquitectónicos) demuestra que el trayecto formativo de nuestra carrera ha rendido sus frutos, cerrando esta etapa universitaria con un broche de oro materializado en este impresionante entregable final.
